\begin{thebibliography}{99}
\bibitem{amour_semiclassical_2013}L. Amour, M. Khodja \& J. Nourrigat - “The semiclassical limit of the time dependent Hartree-Fock equation: the Weyl symbol of the solution”, Anal. PDE 6 (2013) no. 7, p. 1649-1674 | MR | Zbl | DOI
\bibitem{athanassoulis_strong_2011}A. Athanassoulis, T. Paul, F. Pezzotti \& M. Pulvirenti - “Strong semiclassical approximation of Wigner functions for the Hartree dynamics”, Atti Accad. Naz. Lincei Cl. Sci. Fis. Mat. Natur. Rend. Lincei (9) Mat. Appl. 22 (2011) no. 4, p. 525-552 | DOI | MR | Zbl
\bibitem{benedikter_hartree_2016}N. Benedikter, M. Porta, C. Saffirio \& B. Schlein - “From the Hartree dynamics to the Vlasov equation”, Arch. Rational Mech. Anal. 221 (2016) no. 1, p. 273-334 | DOI | MR | Zbl
\bibitem{benedikter_mean-field_2016-1}N. Benedikter, M. Porta \& B. Schlein - “Mean-field regime for fermionic systems”, in Effective evolution equations from quantum dynamics, Springer, Cham, 2016, p. 57-78 | DOI
\bibitem{bergh_interpolation_1976}J. Bergh \& J. Löfström - Interpolation spaces. An introduction, Grundlehren Math. Wiss., vol. 223, Springer, Berlin, Heidelberg, 1976 | DOI
\bibitem{castella_$l^2$_1997}F. Castella - “$L^2$- solutions to the Schrödinger-Poisson system: existence, uniqueness, time behaviour, and smoothing effects”, Math. Models Methods Appl. Sci. 7 (1997) no. 8, p. 1051-1083 | DOI | Zbl
\bibitem{chen_combined_2021}L. Chen, J. Lee \& M. Liew - “Combined mean-field and semiclassical limits of large fermionic systems”, J. Statist. Phys. 182 (2021) no. 2, p. 24 | DOI | MR | Zbl
\bibitem{chong_many-body_2021}J. J. Chong, L. Lafleche \& C. Saffirio - “From many-body quantum dynamics to the Hartree-Fock and Vlasov equations with singular potentials”, 2021 | arXiv
\bibitem{chong_global--time_2022}J. J. Chong, L. Lafleche \& C. Saffirio - “Global-in-time semiclassical regularity for the Hartree-Fock equation”, J. Math. Phys. 63 (2022) no. 8, article ID 081904, 9 pages | DOI | MR | Zbl
\bibitem{figalli_semiclassical_2012}A. Figalli, M. Ligabò \& T. Paul - “Semiclassical limit for mixed states with singular and rough potentials”, Indiana Univ. Math. J. 61 (2012) no. 1, p. 193-222 | MR | Zbl | DOI
\bibitem{golse_schrodinger_2017}F. Golse \& T. Paul - “The Schrödinger equation in the mean-field and semiclassical regime”, Arch. Rational Mech. Anal. 223 (2017) no. 1, p. 57-94 | Zbl | DOI
\bibitem{golse_mean-field_2021}F. Golse \& T. Paul - “Mean-field and classical limit for the $N$-body quantum dynamics with Coulomb interaction”, Comm. Pure Appl. Math. (2021), p. 1-35 | DOI
\bibitem{graffi_mean-field_2003}S. Graffi, A. Martinez \& M. Pulvirenti - “Mean-field approximation of quantum systems and classical limit”, Math. Models Methods Appl. Sci. 13 (2003) no. 1, p. 59-73 | MR | DOI
\bibitem{hunt_lpq_1966}R. A. Hunt - “On $L(p,q)$ spaces”, Enseign. Math. 12 (1966), p. 249-276 | MR | Zbl | DOI
\bibitem{lafleche_propagation_2019}L. Lafleche - “Propagation of moments and semiclassical limit from Hartree to Vlasov equation”, J. Statist. Phys. 177 (2019) no. 1, p. 20-60 | MR | Zbl | DOI
\bibitem{lafleche_global_2021}L. Lafleche - “Global semiclassical limit from Hartree to Vlasov equation for concentrated initial data”, Ann. Inst. H. Poincaré C Anal. Non Linéaire 38 (2021) no. 6, p. 1739-1762 | MR | Zbl | DOI
\bibitem{lafleche_strong_2021}L. Lafleche \& C. Saffirio - “Strong semiclassical limit from Hartree and Hartree-Fock to Vlasov-Poisson equation”, Anal. PDE (2021), to appear
\bibitem{lerner_fefferman-phong_2007}N. Lerner \& Y. Morimoto - “On the Fefferman-Phong inequality and a Wiener-type algebra of pseudodifferential operators”, Publ. RIMS, Kyoto Univ. 43 (2007) no. 2, p. 329-371 | MR | Zbl | DOI
\bibitem{lewin_hartree_2020}M. Lewin \& J. Sabin - “The Hartree and Vlasov equations at positive density”, Comm. Partial Differential Equations 45 (2020) no. 12, p. 1702-1754 | MR | Zbl | DOI
\bibitem{lions_sur_1993}P.-L. Lions \& T. Paul - “Sur les mesures de Wigner”, Rev. Mat. Iberoamericana 9 (1993) no. 3, p. 553-618 | DOI | MR | Zbl
\bibitem{lions_propagation_1991}P.-L. Lions \& B. Perthame - “Propagation of moments and regularity for the 3-dimensional Vlasov-Poisson system”, Invent. Math. 105 (1991) no. 2, p. 415-430 | DOI | MR
\bibitem{markowich_classical_1993}P. A. Markowich \& N. J. Mauser - “The classical limit of a self-consistent quantum Vlasov equation”, Math. Models Methods Appl. Sci. 3 (1993) no. 1, p. 109-124 | DOI | MR | Zbl
\bibitem{narnhofer_vlasov_1981}H. Narnhofer \& G. L. Sewell - “Vlasov hydrodynamics of a quantum mechanical model”, Comm. Math. Phys. 79 (1981) no. 1, p. 9-24 | DOI | MR
\bibitem{pfaffelmoser_global_1992}K. Pfaffelmoser - “Global classical solutions of the Vlasov-Poisson system in three dimensions for general initial data”, J. Differential Equations 95 (1992) no. 2, p. 281-303 | DOI | MR | Zbl
\bibitem{powers_free_1970}R. T. Powers \& E. Størmer - “Free states of the canonical anticommutation relations”, Comm. Math. Phys. 16 (1970) no. 1, p. 1-33 | DOI | MR | Zbl
\bibitem{saffirio_semiclassical_2019}C. Saffirio - “Semiclassical limit to the Vlasov equation with inverse power law potentials”, Comm. Math. Phys. 373 (2019) no. 2, p. 571-619 | DOI | MR | Zbl
\bibitem{saffirio_hartree_2020}C. Saffirio - “From the Hartree equation to the Vlasov-Poisson system: strong convergence for a class of mixed states”, SIAM J. Math. Anal. 52 (2020) no. 6, p. 5533-5553 | DOI | MR | Zbl
\bibitem{spohn_vlasov_1981}H. Spohn - “On the Vlasov hierarchy”, Math. Methods Appl. Sci. 3 (1981) no. 1, p. 445-455 | DOI | MR | Zbl
\end{thebibliography}
